\section{Repair and Refinement Gaps}
\label{sec:preliminaries}

The main theory diagnoses which part of a runtime error is being observed. The representation map $I_t$ determines the information $\calF_t=\sigma(I_t)$ available to the learner. The predictor $f_t$ is the current decision rule. A candidate refinement $\calG\supseteq\calF_t$ is a richer information structure. The two gaps $\Delta_{\mathrm{repair}}$ and $\Delta_{\mathrm{refine}}$ compare these three objects.

\paragraph{Notation.}

All probability spaces are denoted by $(\Omega,\calE,\mathbb P)$ unless the finite case is stated explicitly. Every information structure, including $\calF_t$ and $\calG$, is a sub-sigma-algebra of $\calE$. Random variables are functions on $\Omega$. We reserve $\mathsf A$ for the action space, $I_t$ for the information map, $\calZ_t$ for its label space, $Y$ for the consequence to be predicted, and $U_t$ for newly observed evidence used in a refinement. For a sub-sigma-algebra $\calF$, define
\begin{align*}
\Pred(\calF)=\{a:\Omega\to\mathsf A: a\text{ is }\calF\text{-measurable}\}.
\end{align*}
In the DGM sections, $E_t$ or $E(\calS)$ denotes a graph edge set, $\operatorname{Inc}_t(j)$ denotes the edges incident to concept $j$, and $e_y$ denotes the output embedding of label $y$. We use standard measure-theoretic probability notation for generated sigma-algebras, joins, measurability, and conditional expectation \citep{billingsley1995probability,kallenberg2002foundations}. For vectors $x\in\R^d$, lengths use the $\ell_2$ norm $\norm{x}_2$. For matrices, $\norm{M}_F$ is the Frobenius norm. Inner products are written $\ip{x}{y}$.

\subsection{Information maps and fibers}

At time $t$, let
\begin{align*}
I_t:\Omega\to\calZ_t
\end{align*}
be the learner's current information map. The value $I_t(\omega)$ is the runtime state label available to the learner. Two worlds $\omega,\omega'$ are indistinguishable under the current representation when $I_t(\omega)=I_t(\omega')$. The fiber, or cell,
\begin{align*}
I_t^{-1}(z)=\{\omega:I_t(\omega)=z\}
\end{align*}
is the set of worlds currently treated as the same situation. The sigma-algebra generated by the current information is
\begin{align*}
\calF_t=\sigma(I_t).
\end{align*}

In the finite setting, this is equivalent to a partition
\begin{align*}
\calP=\{B_1,\ldots,B_n\}
\end{align*}
of $\Omega$ into nonempty disjoint cells. A partition $\calQ$ refines $\calP$ if every cell of $\calQ$ is contained in a cell of $\calP$. An event $D\subseteq\Omega$ is $\calP$-measurable if it is a union of cells of $\calP$. Thus a partition is not only notation. It is the learner's current quotient of the world into distinguishable cases.

For sigma-algebras $\calF,\calG\subseteq\calE$, the inclusion $\calF\subseteq\calG$ means that $\calG$ is at least as informative as $\calF$. The join $\calF\vee\calG$ is the smallest sigma-algebra containing both. For a random variable $U$, $\sigma(U)$ is the sigma-algebra generated by $U$; for an event $D$, $\sigma(D)=\{\varnothing,D,D^c,\Omega\}$.

\subsection{Measurable predictors}

A measurable predictor is what the learner is allowed to compute from its current distinctions. With information map $I_t$, the runtime predictor has the form
\begin{align*}
f_t(\omega)=g_t(I_t(\omega))
\end{align*}
for some downstream rule $g_t$. Equivalently, $f_t$ is $\calF_t$-measurable. In a finite partition, a prediction or action $f:\Omega\to\mathsf A$ is $\calP$-measurable if it is constant on each cell:
\begin{align*}
\omega,\omega'\in B\in\calP\quad \Longrightarrow\quad f(\omega)=f(\omega').
\end{align*}
Equivalently, $f$ can depend only on the cell containing the current world state. In the sigma-algebra setting, the admissible predictors are exactly $\Pred(\calF)$. This is the formal reason a missing distinction cannot be repaired by a more clever local update inside the same cell: the update still returns a function on the same quotient.

\subsection{Bayes risk under information}

Bayes risk under information is the best possible performance given the distinctions currently available. Let $Y:\Omega\to\calY$ be a target random variable and let $L:\calY\times\mathsf A\to\R$ be a loss for which the expectations below are well-defined. The Bayes risk available under information $\calF$ is
\begin{align}
R_L(\calF)=\inf_{a\in\Pred(\calF)}\E[L(Y,a)].
\label{eq:bayes-risk}
\end{align}
Thus $R_L(\calF)$ is the best population risk achievable by a predictor that can use only distinctions in $\calF$, following the standard decision-theoretic view of actions, losses, Bayes risk, and comparison of information structures \citep{blackwell1953equivalent,berger1985statistical}.

\subsection{Repair and refinement}

For a concrete predictor $f$, write its risk as
\begin{align*}
R_L(f)=\E[L(Y,f)].
\end{align*}
A \emph{repair} update preserves information:
\begin{align*}
\calF_{t+1}=\calF_t.
\end{align*}
At the population level, a repair can only move within the class of $\calF_t$-measurable predictors, whose optimal risk is $R_L(\calF_t)$. The corresponding repair gap is
\begin{align}
\Delta_{\mathrm{repair}}(f_t,\calF_t)
=
R_L(f_t)-R_L(\calF_t).
\label{eq:repair-gap}
\end{align}

A \emph{refinement} update changes information:
\begin{align*}
\calG\supsetneq\calF_t.
\end{align*}
For a candidate refinement $\calG$, the refinement gap is
\begin{align}
\Delta_{\mathrm{refine}}(\calF_t,\calG)
=
R_L(\calF_t)-R_L(\calG).
\label{eq:refinement-gap}
\end{align}
We omit the loss $L$ from the $\Delta$ notation because the surrounding risk functional already fixes it. The quantity $\Delta_{\mathrm{repair}}$ is residual error under the current information. The quantity $\Delta_{\mathrm{refine}}$ is the value of distinctions not yet present in the current information.

\begin{theorem}[Runtime error decomposition]
\label{thm:error-decomposition}
Let $f_t$ be $\calF_t$-measurable and let $\calG\supseteq\calF_t$ be a candidate refinement. Then
\begin{align}
R_L(f_t)-R_L(\calG)
=
\Delta_{\mathrm{repair}}(f_t,\calF_t)
+
\Delta_{\mathrm{refine}}(\calF_t,\calG).
\label{eq:error-decomposition}
\end{align}
Both terms are nonnegative.
\end{theorem}

\begin{proof}
Add and subtract $R_L(\calF_t)$. The term $\Delta_{\mathrm{repair}}$ is nonnegative because $R_L(\calF_t)$ is the infimum over all $\calF_t$-measurable predictors, including $f_t$. The term $\Delta_{\mathrm{refine}}$ is nonnegative because every $\calF_t$-measurable predictor is also $\calG$-measurable when $\calF_t\subseteq\calG$.
\end{proof}

\subsection{Contradiction inside a cell}

The hard-label obstruction is the clearest special case of positive $\Delta_{\mathrm{refine}}$. If a cell contains states that require different outputs, the current information structure does not merely have a bad parameter value; it lacks the expressive degree of freedom needed to represent the target.

\begin{corollary}[Contradiction requires distinction]
\label{thm:contradiction:informal}
Let $I:\Omega\to\calZ$ be an information map and let $y:\Omega\to\mathsf A$ be a required output. Suppose $I(\omega)=I(\omega')$ but $y(\omega)\ne y(\omega')$. Then no predictor of the form $f=g\circ I$ satisfies both $f(\omega)=y(\omega)$ and $f(\omega')=y(\omega')$.
\end{corollary}

\begin{proof}
If $f=g\circ I$ and $I(\omega)=I(\omega')$, then $f(\omega)=f(\omega')$. Satisfying both required outputs would imply $y(\omega)=y(\omega')$, a contradiction.
\end{proof}

This corollary is deliberately architecture-independent. It applies to a neural cache, a local adapter, a fast-weight memory, a nearest-neighbor system, or a test-time optimizer whenever the update is constrained to the current information. If the relevant states are not separated by $I$, no repair rule inside the cell can represent the required behavior. The contradiction is evidence that the cell should be split.

\begin{example}[One prefix, two consequences]
Suppose the current runtime concept groups two prefixes because their surface form is the same. In one hidden regime the next token should be ``red''; in another it should be ``blue''. Any memory attached to the single concept must store one consequence distribution for both regimes. A repair update can move that distribution, but it cannot make it simultaneously put all mass on ``red'' and all mass on ``blue''. The missing operation is to split the concept by a distinction correlated with the hidden regime, then repair the two resulting memories separately.
\end{example}
